\chapter{解析数论：用连续工具研究离散的素数}

\section{一场越报越吃力的报数游戏}

课间，四个同学玩一个报数游戏：第一个人说“$2$”，后面的人依次报出下一个素数——$3$、$5$、$7$、$11$、$13$……开头大家抢着答，报到五十附近开始要掰手指，报到一百附近就有人掏出草稿纸试除了。游戏有个奇怪的手感：起步时素数密密麻麻，几乎每隔一两个数就撞上一个；越往后，素数之间的“空地”越宽，等待的时间越长。

另一个场景发生在你看不见的地方。你每次用手机支付，浏览器和银行服务器之间都要交换密钥，而最流行的公开密钥方案之一，第一步就是“找两个几百位长的大素数”。计算机怎么找？它不是背一张素数表——几百位的素数全列出来，全宇宙的原子都不够用。它的办法笨得可爱：随机取一个大数，测试它是不是素数，不是就换一个，直到撞上一个为止。这个办法能行，全靠一个背景事实：即使到了几百位那么高的地方，素数虽然稀，却没有稀到绝迹，平均每试几百次总能中一次。

于是两个问题浮上来。第一个问题我们其实已经解决过：素数会不会用完？不会，欧几里得两千多年前就证明了素数有无穷多个。但第二个问题要刁钻得多，也正是这场报数游戏“越报越吃力”的原因：

\begin{quote}
素数越到后面越稀。可是——稀到什么程度？稀得有没有规律？能不能写出一个公式，大致告诉我们：在数字 $x$ 附近，素数出现的“密度”是多少？
\end{quote}

这个问题听起来像统计，像物理，不太像我们习惯的“整数问题”。整数是一个一个蹦的，是离散的；而“密度”“规律”“公式”这些词，属于连续变化的世界。本章要讲的解析数论，干的正是这件看起来不讲道理的事：用连续的数学工具——对数、积分、无穷级数，甚至复数——去研究最离散的素数。令人吃惊的是，这条路不但走得通，还引出了数学中最著名的未解之谜。

\section{第一个猜想：按比例分钱为什么行不通}

先凭直觉猜。一百以内有 $25$ 个素数，占四分之一。于是最自然的猜想是：素数大约占整数的四分之一，往后也大致如此，只是慢慢再少一点。换句话说，素数的密度近似一个固定比例。

这个猜想只要多算一段就破产了。如果你耐心地把素数表数下去（或者查一张现成的表），会发现：$100$ 以内有 $25$ 个素数，密度是 $25\%$；$1000$ 以内有 $168$ 个，密度降到 $16.8\%$；$10000$ 以内有 $1229$ 个，密度只剩 $12.3\%$。密度不是在一个固定值附近波动，而是持续、稳定地往下掉。任何“固定比例”的猜法都会被这张表否决。

第二个直觉是：密度会不会像 $\dfrac{1}{x}$ 那样下降？也就是到 $x$ 附近时，大约每 $x$ 个数里才有 $1$ 个素数？这也错得太快——按这个说法，$1000$ 附近每 $1000$ 个数才有 $1$ 个素数，可实际上 $1000$ 以内有 $168$ 个素数，平均不到 $6$ 个数就有 $1$ 个。真实的下降速度比 $\dfrac{1}{x}$ 慢得多，又比“固定比例”快得多。它卡在中间，不上不下。

\begin{fanli}
“素数越来越稀，所以到了足够大的地方，就会有一段很长很长的区域一个素数都没有，长得超过任何指定的长度——这说明素数‘实际上是有限的’。”

前半句是对的（我们很快会证明：任意长的纯合数段确实存在），后半句是错的。稀，不等于消失。欧几里得的证明保证：无论你走多远，前面永远还有素数。直觉上“稀到看不见”和逻辑上“没有了”是两回事，这正是本章要反复小心的地方。
\end{fanli}

固定比例失败，$\dfrac{1}{x}$ 也失败。我们需要一个会“缓慢衰减”的新函数来当密度的候选人。这个函数其实你在第三部已经见过——对数。

\section{给素数造一台计数器}

要谈“密度”，先要有一把尺子。我们发明的第一个新工具是一个计数函数：

\begin{dingyi}[素数计数函数]
对任意正实数 $x$，用 $\pi(x)$ 表示不超过 $x$ 的素数个数。例如 $\pi(10)=4$（素数是 $2,3,5,7$），$\pi(100)=25$。
\end{dingyi}

这里的 $\pi$ 不是圆周率，只是借用了希腊字母“p”——英语里素数是 prime。这个符号解决的麻烦是：以前我们只能说“素数好像越来越少”，现在我们可以说“$\pi(100)=25$，$\pi(1000)=168$”，把模糊的感觉变成可以比较、可以作图的数。

$\pi(x)$ 的图像是什么样子？它是一个一个台阶：$x$ 走过合数时它原地不动，每踩到一个素数就往上跳一格。它永远向右走，永远只上不下，而且（因为素数无穷多）台阶会一直延伸下去，永不到顶。

\begin{center}
\begin{tikzpicture}[scale=1.0]
  \draw[->] (0,0) -- (6.5,0) node[below left] {$x$};
  \draw[->] (0,0) -- (0,4.2) node[below left] {$\pi(x)$};
  \draw[thick,blue!60!black]
    (0,0) -- (0.4,0) -- (0.4,0.35) -- (0.6,0.35) -- (0.6,0.7)
    -- (1.0,0.7) -- (1.0,1.05) -- (1.4,1.05) -- (1.4,1.4)
    -- (2.2,1.4) -- (2.2,1.75) -- (2.6,1.75) -- (2.6,2.1)
    -- (3.4,2.1) -- (3.4,2.45) -- (3.8,2.45) -- (3.8,2.8)
    -- (4.6,2.8) -- (4.6,3.15) -- (5.8,3.15) -- (5.8,3.5) -- (6.3,3.5);
  \foreach \p/\v in {0.4/2, 1.0/5, 2.2/11, 3.4/17, 4.6/23, 5.8/29}
    \node[below] at (\p,0) {\scriptsize $\v$};
  \node[right] at (6.3,3.5) {\scriptsize 每跳一格，就是一个素数};
\end{tikzpicture}
\end{center}

这台“计数器”是纯离散的：台阶高矮错落，看不出光滑的规律。真正神奇的事情是——如果站远一点看，这条坑坑洼洼的楼梯，竟然贴着一条光滑的曲线走。为了说清楚这条曲线，我们需要第二位主角：自然对数 $\ln x$。

回忆一下对数的出身。指数函数 $10^{2}=100$、$10^{3}=1000$ 告诉我们：$\lg 100=2$、$\lg 1000=3$。对数做的是“把乘法翻译成加法”的买卖，它增长得极慢：$x$ 翻十倍，$\lg x$ 只加一。数学家更偏爱以 $e\approx2.718$ 为底的自然对数 $\ln x$，因为它在微积分里最顺手：曲线 $y=\dfrac{1}{t}$ 从 $1$ 到 $x$ 下方的面积，恰好就是 $\ln x$。这个“面积”身份，待会儿要派大用场。

第三个要请出场的工具，是一条名声很大的级数：

\[
1+\frac12+\frac13+\frac14+\frac15+\cdots
\]

它叫调和级数。每一项都在变小，趋向于零，可它加起来的总和却是无穷大！看出这一点的办法很巧：把项分组——$\frac12$ 单独一组；$\frac13+\frac14$ 一组，两项都不小于 $\frac14$，和至少是 $\frac12$；接下来的 $\frac15$ 到 $\frac18$ 共四项，每项不小于 $\frac18$，和又至少是 $\frac12$；再往后八项、十六项……每一组的和都至少 $\frac12$。无穷多个 $\frac12$ 加起来，当然发散到无穷大。调和级数给我们的教训是：趋向于零，不等于加起来有限。

调和级数还有一个关键的量化事实：前 $n$ 项的和，大约等于 $\ln n$。这不奇怪——上面说过，$\ln n$ 就是曲线 $y=\frac1t$ 下方的面积，而 $1+\frac12+\cdots+\frac1n$ 不过是把这块面积切成宽度为 $1$、高度为 $\frac1k$ 的小矩形之和，矩形和面积本来就差不了多少。粗略地说：

\[
1+\frac12+\frac13+\cdots+\frac1n \;\approx\; \ln n.
\]

现在，三件工具都摆上桌了：台阶状的计数器 $\pi(x)$、慢吞吞的对数 $\ln x$、以及对数与“倒数之和”之间的面积桥梁。是时候说出本章的中心故事了。

\section{素数定理：稀疏中的规律}

十八世纪末，少年高斯得到一个好玩的机会：有人送他一本对数表，书的附录里印着一张素数表。据说他翻来覆去地比对这两张表，注意到了一件怪事——素数出现的“疏密节奏”，和对数增长的节奏，竟然合拍。他猜：在 $x$ 附近，每 $\ln x$ 个连续整数里，大约有 $1$ 个素数。换句话说，$x$ 附近的素数密度大约是 $\dfrac{1}{\ln x}$。

密度是 $\dfrac{1}{\ln x}$，那么不超过 $x$ 的素数总数，就该是把这个密度从 $2$ 一路累加到 $x$——用积分的语言写，这条光滑曲线就是 $\displaystyle\int_2^x \frac{\mathrm{d}t}{\ln t}$。一个更粗糙但更好算的近似是 $\dfrac{x}{\ln x}$。高斯的猜测后来被证明是对的，这就是素数定理：

\begin{dingli}[素数定理]
当 $x$ 越来越大时，
\[
\pi(x)\;\approx\;\frac{x}{\ln x},
\]
精确的意思是：比值 $\pi(x)\big/\dfrac{x}{\ln x}$ 随着 $x$ 无限增大而趋向于 $1$。
\end{dingli}

\begin{zhu}
高斯在十八世纪末猜出这条规律，但证明等了一百年。直到 $1896$ 年，数学家阿达马和德拉瓦莱普桑才各自独立地证明了它，而他们用的武器不是整数技巧，恰恰是复数平面上的函数理论——也就是本章后半部分的主角。素数定理的证明超出了本书范围，但它的内容我们可以用表格亲手检验。
\end{zhu}

动手验一验。$\ln 100\approx4.6$，$\ln 1000\approx6.9$，$\ln 10^{6}\approx13.8$，$\ln 10^{9}\approx20.7$。填进 $\dfrac{x}{\ln x}$，和真实的 $\pi(x)$ 对照：

\begin{center}
\begin{tabular}{crrc}
\toprule
$x$ & $\pi(x)$ & $x/\ln x$（取整） & 比值 $\pi(x)/(x/\ln x)$ \\
\midrule
$10^{2}$ & $25$ & $22$ & $1.15$ \\
$10^{3}$ & $168$ & $145$ & $1.16$ \\
$10^{4}$ & $1229$ & $1086$ & $1.13$ \\
$10^{5}$ & $9592$ & $8686$ & $1.10$ \\
$10^{6}$ & $78498$ & $72382$ & $1.08$ \\
$10^{9}$ & $50847534$ & $48254942$ & $1.05$ \\
\bottomrule
\end{tabular}
\end{center}

看最右一列：$1.15,1.16,1.13,1.10,1.08,1.05$，慢吞吞地、一点一点地朝 $1$ 蹭过去。这就是“比值趋向于 $1$”在现实中的样子——不急，但从不停下，也从不回头。一张表当然不构成证明（这一点我们后面还要强调），但它让“素数的稀疏有精确规律”从一句口号变成看得见的数。

把素数定理翻译回报数游戏：在 $x$ 附近，素数的平均间隔大约是 $\ln x$。$100$ 附近平均每隔约 $4.6$ 个数有一个素数；$10^{6}$ 附近平均每隔约 $13.8$ 个数才有一个；到几百位的天文数字那里，平均每隔几百个数有一个——这正好解释了为什么电脑“随机试几百次就能撞上一个大素数”。密码学家不是运气好，是素数定理在背后托底。

\begin{shuyu}{素数定理}
\textbf{直观解释：} 素数越来越稀，但稀疏的节奏被 $\ln x$ 精确控制；$x$ 附近的素数密度约为 $1/\ln x$。\\
\textbf{正式表述：} $\pi(x)/(x/\ln x)\to 1$（当 $x$ 无限增大）。\\
\textbf{一个例子：} $x=10^{6}$ 时，$x/\ln x\approx72382$，真实的 $\pi(x)=78498$，两者只差约 $8\%$，而且 $x$ 越大这个百分比越小。
\end{shuyu}

不过要小心：“平均间隔约 $\ln x$”说的是平均，不是每一段都如此。素数的真实间隔忽大忽小：有时两个素数只差 $2$，比如 $29$ 和 $31$、$41$ 和 $43$，这叫孪生素数；有时又一连成百上千个合数。事实上，我们能证明“空地”可以任意长：

\begin{dingli}[素数间隔可以任意大]
对任意正整数 $n$，都存在连续 $n$ 个正整数，其中没有一个是素数。
\end{dingli}

\begin{proof}
取一个足够大的数 $N=(n+1)!$，即 $1\times2\times3\times\cdots\times(n+1)$ 的乘积。考察紧跟在 $N$ 后面的 $n$ 个数：
\[
N+2,\ N+3,\ N+4,\ \dots,\ N+(n+1).
\]
第一个数 $N+2$：$N$ 本身含有因子 $2$（阶乘把 $2$ 乘进去了），所以 $N+2$ 能被 $2$ 整除，是合数。第二个数 $N+3$：$N$ 含有因子 $3$，所以 $N+3$ 能被 $3$ 整除，也是合数。一般地，对 $2$ 到 $n+1$ 之间的每个 $k$，$N$ 都含有因子 $k$，于是 $N+k$ 能被 $k$ 整除，而且 $N+k$ 显然大于 $k$，所以它是合数。这样，$N+2$ 到 $N+(n+1)$ 这 $n$ 个连续整数全是合数。
\end{proof}

这个证明用的是纯构造：想要多长的空地，就把阶乘造多大。比如想要连续 $1000$ 个合数，取 $1001!$ 加 $2$ 到 $1001!$ 加 $1001$ 即可——一个天文数字长度的“素数沙漠”。可是另一头，孪生素数似乎又永远成双成对地冒出来：$3$ 和 $5$、$11$ 和 $13$、$\dots$、$101$ 和 $103$、$\dots$。数学家猜孪生素数有无穷多对——这是著名的孪生素数猜想，至今没人能证明，也没人能推翻。$2013$ 年，张益唐证明了存在无穷多对间隔不超过七千万的素数，随后全世界的合作把这个上界一路压缩到 $246$。从 $246$ 到 $2$，看似一步之遥，至今无人跨过。

于是素数间隔的图景是：平均而言越来越宽（受 $\ln x$ 控制），个别地方却任意宽，而相隔 $2$ 的“双胞胎”又似乎永不消失——平均的规律性，和个体的任性，同时成立。要搞清楚“素数究竟有多规律”，光有 $\pi(x)\approx x/\ln x$ 还不够，我们还得问：误差有多大？$\pi(x)$ 和 $\dfrac{x}{\ln x}$ 之间的差，是乱蹦的，还是被什么更深的东西管着？这个问题，把我们带到了数学最神秘的大门口。

\section{一台把素数“调出来”的级数机器}

让我们把镜头从表格转向一台“机器”。十八世纪，欧拉研究下面这个无穷级数：

\[
\zeta(s)\;=\;1+\frac{1}{2^{s}}+\frac{1}{3^{s}}+\frac{1}{4^{s}}+\cdots
\]

这里的 $s$ 先当作大于 $1$ 的实数。$s=1$ 时它就是刚才的调和级数，发散；$s=2$ 时它收敛到 $\dfrac{\pi^{2}}{6}$（这个结果本身就是一段传奇，圆周率竟然从整数的平方倒数里冒出来）。欧拉注意到的关键事实是：这台看似与素数无关的级数机器，骨子里完全由素数决定。原因是算术基本定理——每个整数都唯一地分解成素数的乘积。把级数里每一项的分母做素数分解、重新归拢，利用等比数列求和，就能把加法改写成乘法：

\[
1+\frac{1}{2^{s}}+\frac{1}{3^{s}}+\frac{1}{4^{s}}+\cdots
\;=\;
\frac{1}{1-2^{-s}}\times\frac{1}{1-3^{-s}}\times\frac{1}{1-5^{-s}}\times\frac{1}{1-7^{-s}}\times\cdots
\]

右边是一个跑遍所有素数的无穷乘积。这个等式叫欧拉乘积公式，它是本章最深的一句诗：左边是所有整数的合唱，右边是每个素数各管一件乐器。整数的加法世界和素数的乘法世界，在这台机器里是同一件东西的两种写法。

立刻有一个推论。让 $s$ 趋向于 $1$，左边的级数发散到无穷大；如果素数只有有限个，右边只是有限个分数相乘，必定是有限的，矛盾。所以素数必须无穷多——欧拉乘积给了“素数无穷”一个全新风格的证明。更精细地算下去，还能得到：所有素数的倒数之和

\[
\frac12+\frac13+\frac15+\frac17+\frac1{11}+\cdots
\]

也是发散的，而且它发散的速度大约像 $\ln(\ln x)$ 那样慢。这个数字值得咂摸：素数在全体整数中的“分量”，重到倒数和都收不住——比起平方数（$\sum\frac{1}{n^2}$ 收敛），素数要“稠密”得多。素数无穷，不是勉强无穷，而是理直气壮地无穷。

\section{黎曼的惊天一跃：把 $s$ 送进复数平面}

欧拉的机器要求 $s>1$，否则级数收不住。$1859$ 年，黎曼只发表了一篇八页的短文，却做了一件改变数学史的事：他问，能不能把 $\zeta(s)$ 的定义域扩展到整个复数平面？

\begin{shuyu}{黎曼 $\zeta$ 函数}
\textbf{直观解释：} 一台以复数为输入的机器，$s>1$ 时由级数 $1+2^{-s}+3^{-s}+\cdots$ 驱动；通过“解析延拓”的技术，它几乎可以延伸到整个复数平面。\\
\textbf{正式定义（的部分）：} $\zeta(s)=\sum_{n=1}^{\infty} n^{-s}$（当 $s$ 的实部大于 $1$），再唯一地延拓到其余地方，仅在 $s=1$ 处“爆炸”。\\
\textbf{一个例子：} $\zeta(2)=\dfrac{\pi^{2}}{6}$；而在延拓意义下 $\zeta(-1)$ 被赋予值 $-\dfrac{1}{12}$——注意这不是说 $1+2+3+\cdots$ 真的等于 $-\frac{1}{12}$，而是延拓后的函数在那里的取值，两者不能混为一谈。
\end{shuyu}

“延拓”这一步的难处在于：级数本身只在实部大于 $1$ 的区域收敛，出了这个圈子，$1+2^{-s}+3^{-s}+\cdots$ 根本不听话。数学家用了一套精巧的办法（大意是先绕个弯把级数改写成收敛的样子，再证明这样得到的函数是唯一合理的延伸），让 $\zeta$ 在除 $s=1$ 以外的整个复数平面上都有确定的意义。这套技术叫解析延拓，属于复分析，细节超出本书——但请记住它的一个宝贵性质：延拓是唯一的，就像把一段圆弧顺着自己的弯曲方式继续画下去，只有唯一一种自然的画法。

函数有了，黎曼接着问：这台机器在哪些输入下输出为零？也就是 $\zeta(s)=0$ 的解在哪里？零点是函数的“基因”——就像多项式由它的根决定一样，一个解析函数在很大程度上由它的零点决定。计算表明，负偶数 $-2,-4,-6,\dots$ 都是零点，这些被称为平凡零点，没什么好惊讶的。真正的秘密藏在实部介于 $0$ 和 $1$ 之间的一条竖直“临界带”里：那里还有无穷多个非平凡零点，而且黎曼算出，最前面几个的实部恰好都是 $\dfrac12$：

\[
\frac12+14.13\,\mathrm{i},\qquad
\frac12+21.02\,\mathrm{i},\qquad
\frac12+25.01\,\mathrm{i},\qquad\dots
\]

\begin{center}
\begin{tikzpicture}[scale=1.0]
  \draw[->] (-0.5,0) -- (2.5,0) node[below left] {实部};
  \draw[->] (0,-0.4) -- (0,3.8) node[below left] {虚部};
  \draw[dashed,gray] (1,-0.3) -- (1,3.6);
  \node[below,gray] at (1,0) {\scriptsize $\frac12$};
  \node[below] at (2,0) {\scriptsize $1$};
  \node[above right,gray] at (1,3.5) {\scriptsize 临界直线};
  \foreach \y in {1.7, 2.52, 3.0}
    \fill[red!60!black] (1,\y) circle (1.6pt);
  \node[right] at (1.1,1.7) {\scriptsize $\frac12+14.13\mathrm{i}$};
  \node[right] at (1.1,2.52) {\scriptsize $\frac12+21.02\mathrm{i}$};
  \node[right] at (1.1,3.0) {\scriptsize $\frac12+25.01\mathrm{i}$};
  \foreach \y in {0.85, 2.1, 2.75}
    \fill[blue!50!black] (0.4,\y) circle (1.2pt);
  \node[left,blue!50!black] at (0.3,0.85) {\scriptsize 猜想的“如果”};
\end{tikzpicture}
\end{center}

为什么零点与素数有关？因为 $\zeta$ 函数通过欧拉乘积“装着”全部素数的信息，而零点就是提取这些信息的旋钮。黎曼在那篇八页论文里给出了一个惊人的公式（证明后来被补全）：$\pi(x)$ 可以写成一项光滑的主项，加上每个非平凡零点贡献的一项“修正波”。零点的实部越大，它贡献的波浪起伏越厉害，$\pi(x)$ 偏离光滑曲线的误差就越难控制。于是“素数有多规律”这个问题，被精确地翻译成了“零点躺在哪里”这个几何问题。

黎曼把他算到的现象写成了一个猜想：

\begin{quote}
\textbf{黎曼猜想（1859 年，至今未证明）：} $\zeta$ 函数的所有非平凡零点，实部都等于 $\dfrac12$——它们整整齐齐地躺在同一条竖直直线上。
\end{quote}

请注意措辞：这是\textbf{猜想}，不是定理。一百六十多年里，计算机验证了超过十万亿个非平凡零点，无一例外全在这条线上；无数数学家尝试证明它，相关的等价命题和推论堆成了山，但证明至今没有出现。它是克雷数学研究所悬赏一百万美元的七个“千禧年难题”之一。

黎曼猜想为什么重要？因为它管着素数定理的误差。数学家已经证明：$\pi(x)$ 与光滑曲线之间的偏差大小，完全由非平凡零点的实部控制。如果黎曼猜想成立，这个偏差就被压到最小，大致不超过 $\sqrt{x}\,\ln x$ 的量级——相对于 $x$ 本身，这就像测量地球到月球的距离时误差不到一公里。反过来，只要有一个零点偏离临界线，素数分布就会出现比预期大得多的“异常起伏”。所以黎曼猜想的通俗版本就是一句话：\textbf{素数虽然顽皮，但顽皮得有底线；它们的随机性被一条直线死死框住。}

\begin{fanli}
“计算机已经验证了十万亿个零点都在临界线上，所以黎曼猜想基本算是成立了，只差走个形式。”

这是把实验当成证明的典型错误。非平凡零点有无穷多个，验证再多也只是无穷海洋里的一瓢水。数学史上有过残酷的教训：有的命题对极大极大的数都成立，却在更远的、人力无法企及的地方出现第一个反例（素数分布的研究里就真实发生过“巨大才反转”的现象）。实验的价值在于启发猜想、积累信心、寻找反例；而“对所有情形成立”的保证，只能来自证明。这正是本章标题里“解析”二字的分量：它要求的是论证，不是统计。
\end{fanli}

\begin{lianjie}
回头看，本章把前面好几章的线拧成了一股绳：第 9 章的极限（比值趋向于 $1$）、第 11 章的级数与积分（调和级数、$\ln x$ 的面积身份）、第 12 章的复数（$\zeta$ 函数的延拓）、第 13 章的算术基本定理（欧拉乘积的基石）和第 14 章的计算实验（素数表与筛法）。向前看，黎曼猜想里“所有零点躺在一条直线上”这种关于“形状与位置”的表述，已经在向第 16 章的拓扑学招手。而密码学里的公开密钥方案，则是素数分布在现实世界最硬的落地：它一边依赖“找素数容易”（素数定理保证密度），一边依赖“分解大数困难”（至今没有快速算法）。
\end{lianjie}

\section{亲手数一数：实验能做到什么，做不到什么}

\begin{shiyan}
准备一张 $1$ 到 $100$ 的方格表，亲手做一遍埃拉托色尼筛法：圈出 $2$，划掉它的倍数；圈出下一个幸存的 $3$，划掉它的倍数；再对 $5$、$7$ 各做一遍（$100$ 以内只需筛到 $\sqrt{100}=10$ 即可停手，想一想为什么）。最后幸存的数就是 $100$ 以内的全部素数，数一数，是不是 $25$ 个？

接着记录每相邻两个素数的间隔：$3-2=1$，$5-3=2$，$7-5=2$，$11-7=4$，……把间隔列成一行，回答三个问题：
\begin{enumerate}
  \item 最大的间隔是几？出现在哪里？
  \item 间隔为 $2$ 的孪生素数有几对？
  \item 平均间隔大约是多少？和 $\ln 100\approx 4.6$ 比一比，接近吗？
\end{enumerate}
如果你有计算器或电脑，可以再把实验做到 $1000$，观察“平均间隔”如何向 $\ln 1000\approx 6.9$ 移动。有条件的同学还可以把这段过程写成十几行的短程序：用两重循环模仿上面的“圈与划”，让计算机替你数到十万。
\end{shiyan}

做完实验，请停下来想一想：你手里的表格和程序，到底证明了什么？它告诉你“$100$ 以内素数的平均间隔约为 $4.3$，与 $\ln 100$ 同量级”——这是一次观察，是素数定理在 $x=100$ 处的一次抽样。它没有、也不可能证明比值趋向于 $1$ 的极限事实。同样的逻辑链条也适用于黎曼猜想：十万亿次验证依然是实验。数学里有三层台阶，本章恰好各占了一级：\textbf{实验}（筛法、数表、零点计算）负责提供线索，\textbf{猜想}（孪生素数、黎曼猜想）负责指出方向，\textbf{证明}（欧几里得、阶乘构造、素数定理）负责给出保证。混淆这三层，是数论里最容易掉进去的坑，也是最有资格被记住的教训。

\begin{toolbox}
本章新添的工具：
\begin{itemize}
  \item \textbf{计数函数 $\pi(x)$}：把“素数越来越稀”的感觉变成可作图、可比较的函数。
  \item \textbf{自然对数 $\ln x$}：增长最慢的基本函数之一，恰好是素数密度的“刻度尺”；它同时是 $y=\frac1t$ 下方的面积。
  \item \textbf{调和级数与它的发散}：趋向于零不等于收敛；前 $n$ 项和约为 $\ln n$。
  \item \textbf{欧拉乘积}：$\sum n^{-s}=\prod_{p}(1-p^{-s})^{-1}$，整数加法世界与素数乘法世界的翻译器。
  \item \textbf{$\zeta$ 函数与解析延拓}：把级数机器延伸到复数平面，零点的位置控制素数分布的误差。
  \item \textbf{实验—猜想—证明的分层意识}：表格启发猜想，猜想指引证明，只有证明给出保证。
\end{itemize}
\end{toolbox}

\begin{pandeng}
想继续往前走的同学，可以记住这些路标：\textbf{切比雪夫函数}与素数定理的初等证明（$1949$ 年塞尔伯格与埃尔德什各自给出，不依赖复分析）；\textbf{对数积分} $\mathrm{Li}(x)$，它比 $x/\ln x$ 更贴近 $\pi(x)$；\textbf{解析延拓}与复分析，黎曼论文的真正语言；\textbf{筛法}的现代版本，它们支撑着张益唐之后的素数间隔研究；\textbf{零点计算}，人们如何用高速计算机把十万亿个零点挨个钉在临界线上。这些关键词中的任何一个，都够你在图书馆或网络上泡一个暑假。
\end{pandeng}

\section*{思考题}

\textbf{观察题}
\begin{sikao}
  \item 重做 $100$ 以内素数的间隔记录，找出所有“间隔突然变大”的位置（即比它前面所有间隔都大的那些）。观察这些位置上的素数，猜一猜：下一个纪录间隔大概会出现在什么样的地方？\emph{提示：纪录间隔出现时，意味着一段较长的合数沙漠；想想阶乘构造给你的感觉——沙漠更可能出现在数的“结构特殊”处，还是随便哪里？}
  \item 用正文的表格观察“比值”一列：从 $10^{2}$ 到 $10^{9}$，比值每下降一点，$x$ 要扩大多少倍？由此感受“趋向于 $1$”到底有多慢。\emph{提示：注意 $x$ 每乘以 $10$，比值只挪动零点几个百分点；这不是它不努力，而是 $\ln x$ 增长得太矜持。}
\end{sikao}

\textbf{证明题}
\begin{sikao}
  \item 证明：存在连续 $1000$ 个正整数，其中没有一个素数。并说明你构造出的那段合数大约有多“高”。\emph{提示：把正文里的 $N=(n+1)!$ 换成合适的阶乘；检查每个 $N+k$ 为什么一定是合数。}
  \item 不用查表，证明调和级数中任意取出充分靠后的连续若干项，它们的和可以超过 $10$。\emph{提示：回忆正文的分组办法——每一组的和至少是 $\frac12$；要让总和超过 $10$，大约需要多少个这样的组？}
\end{sikao}

\textbf{挑战题}
\begin{sikao}
  \item 正文用欧拉乘积说明了“若素数只有有限个，则 $s\to1$ 时两边矛盾”。顺着这个思路再进一步：试说明为什么“素数倒数和发散”这件事，从乘积公式的角度看几乎是必然的。\emph{提示：对乘积公式两边取对数，把乘积变成求和；当 $s\to1$ 时左边爆炸，右边那些 $\ln(1-p^{-s})$ 展开后领头的项正是 $\frac1p$。这一步需要你对“取对数把乘法变加法”足够熟练，卡住是正常的。}
\end{sikao}

\section*{通往下一章的门}

这一章我们干了一件“跨界”的事：用连续的、光滑的分析工具，去审问离散的、一粒一粒的素数。回头看那幅零点的示意图，有一个细节值得再盯一眼——黎曼猜想说的是：无穷多个神秘的点，全部躺在\textbf{同一条直线}上。

“躺在一条线上”“围成一个圈”“中间有一个洞”——这类话我们已经随口用了好几次。可是，什么叫一条线？什么叫一个洞？如果把整个复数平面换成一块橡皮膜，允许拉伸、弯折、扭转，只要不撕破、不粘断，那条临界线还是“一条线”吗？咖啡杯和甜甜圈，在只许揉不许撕的规则下，凭什么算是同一种东西，又凭什么和球不是同一种东西？

数论靠计数和证明前进，而回答上面这些问题，需要一门专门研究“揉不坏的性质”的学问。它不关心长度和角度，只关心连不连、通不通、有几个洞。这就是下一章——拓扑学——要带你进入的世界。
